%------------------------------------------------------------------------------%
\chapter{Funções}

Esse capítulo será dedicado a apresentar o conceito de função e apresentar
algumas funções usuais. Esses conceitos são importantes e servirão de base para
o estudo de funções específicas nas aulas posteriores.

%------------------------------------------------------------------------------%
\section{Funções}

As funções surgem quando uma quantidade (Grandeza) depende de outra. Por
exemplo, podemos dizer que a área de um quadrado está relacionada ao
comprimento do lado, isso quer dizer que uma vez conhecido o comprimento do
lado, $x$, conseguimos calcular a área usando a função $f$ descrita pela
relação $f(x)=x^2$.

É possível representar uma função de quatro maneiras:
\begin{itemize}
  \item verbalmente (descrevendo-a com palavras)
  
  \item numericamente (por meio de uma tabela de valores)
  
  \item visualmente (através de um gráfico)
  
  \item algebricamente (utilizando-se uma fórmula explícita)
\end{itemize}

\begin{defin}(\textbf{Função}): 
  Uma \textbf{função} $f:D\to C$ é uma relação que associa a cada elemento $x$
  de um conjunto $D$, chamado \textbf{domínio}, um \textit{único} elemento
  $f(x)$ (ou $y$) de um conjunto $C$, denominado \textbf{contradomínio}.
\end{defin}

\begin{obs} 
  Quando o domínio de uma função $f$ não for específicado, então o
  \textbf{domínio} (natural) de $f$ é o conjunto de todos os valores de $x$
  para os quais $f$ está definida.
\end{obs}

\begin{obs}
  Devemos ter cuidado ao analisarmos o domínio de uma função, quando essa é
  originada de um problema real. 
  Se, por exemplo, a função $f(x)=x+1$, $f(x)$
  representa o valor a ser pago em uma corrida de táxi em função do número de
  quilômetros rodados $x$, então o domínio de $f$ será 
  $D= \{ x \in \R | x \geq 0 \}$, 
  uma vez que o número de quilômetros rodados é sempre não-negativo.
\end{obs}

Note que o contradomínio $C$ de uma função $f$ pode ser qualquer conjunto que
contenha, dentre outros elementos, os valores de $f(x)$. Sendo assim, em muitas
situações estamos interessados em trabalhar apenas com o conjunto que contém
apenas os valores gerados por $f$, ao qual damos o nome de \textbf{conjunto
imagem}.

\begin{defin}(\textbf{Imagem})
  Dada uma \textbf{função} $f:D\to C$, o valor $f(x)$ para $x\in D$ é chamado
  \textbf{imagem de $x$}. O conjunto de todas as imagens da função $f$ é
  denominado \textbf{conjunto imagem}(ou simplesmente imagem) e denotado por
  \textbf{\textit{Im}(f)}. Ou seja,  \textit{Im}(f)$=\{f(x);x \in D\}$.
\end{defin}

\begin{defin} (\textbf{Gráfico}) 
  o \textbf{gráfico} de $f$ é o conjunto dos pontos $(x,y)$ do plano cartesiano
  que satisfazem $y=f(x),$ ou seja,  $G(f)=\{(x, f(x)); x\in D\}$.
\end{defin}

\begin{figure}[h!]
  \includegraphics[scale=0.5]{funcao1}
  \caption{Figura retirada do livro Cálculo 1, James Stewart.}
\end{figure}

\textbf{Quando que uma aplicação não é função?}

\begin{exe}
  Verifique se a equação $y^2+x^2-16=0$ permite uma definição 
  de $y$ como função de $x$.
\end{exe}

Inicialmente vamos ver se conseguimos explicitar $y$ como função de $x$.
\begin{align*}
y^2 & =    16-x^2 \\
y   & =\pm \sqrt{16-x^2}
\end{align*}
Logo, para cada $x \in \mathbb{R}$, temos dois valores para $y$. Portanto, a
equação $y^2+x^2-16=0$ não permite uma definição de $y$ como função de $x$.

\begin{exe}
  Faça o estudo do domínio das seguintes funções \\
  a) $f(x)=1-x^4$                            \hspace{1cm}
  b) $g(x)=\frac{1}{x+4}$                    \hspace{1cm}
  c) $h(x)=\sqrt{9-x}$                       \hspace{1cm}
  d) $f(t)=\frac{\sqrt{2+t}}{\sqrt{2t-6}}$   \hspace{1cm}  
  e) $f(x)=\sqrt[3]{2x-3}$                   \hspace{1cm} 
  f) $p(x)=\frac{\sqrt{x-1}}{x^2-5x+6}$      \hspace{1cm} \\

  \textbf{a)} Note que $f$ está bem definida para todo $x \in \mathbb{R}$, logo
  $D(f)= \mathbb{R}$. \\

  \textbf{b)} Nesse caso temos que $g$ está definida quando $x+4 \neq 0
  \Rightarrow x \neq -4$, logo $D(g)= \mathbb{R}-\{-4\}$ ou $D(g)= \{ x \in
  \mathbb{R}| x \neq -4 \}$. \\
  
  \textbf{c)} Para que $h$ seja bem definida $9-x \geq 0 \Rightarrow  x \leq 9
  $, logo $D(h)= \{ x \in \mathbb{R}| x \leq 9 \}$. \\
  
  \textbf{d)} Para que $f$ esteja bem definida teremos que ter $2+t \geq 0
  \Rightarrow t \geq -2$ e $2t-6 >0 \Rightarrow t>3$, fazendo a interseção
  desses intervalos temos que $D(f)= \{ x \in \mathbb{R}| x \geq 3 \}$. \\
  
  \textbf{e)} Como o índice a raíz é ímpar, então $D(f)= \mathbb{R}$. \\
  
  \textbf{f)} Para que $p$ esteja bem definida teremos que ter 
  $x-1 \geq 0 \Rightarrow x \geq 0$ 
  e 
  $x^2-5x+6 \neq 0 \Rightarrow x \neq 2$ 
  e 
  $x \neq 3$,
  portanto 
  $D(p)= \{ x \in \R | x \geq 1, x \neq 2 \hspace{0,15cm} x \neq 3 \}$.
\end{exe}

%------------------------------------------------------------------------------%
\subsection{Simetria}

Se uma função $f$ satisfaz $f (–x)$ = $f (x)$ para todo número $x$ em seu
domínio, então $f$ é chamada \textbf{função par}. Por outro lado, se $f$
satisfaz $f (–x) = –f (x)$ para cada número $x$ em seu domínio, uma $f$  é
\textbf{função ímpar}.

\begin{exe}
  Verifique se as funções abaixo são par, ímpar ou nem par, nem ímpar \\
  a) $f(x)=1-x^4$  \hspace{1cm}   
  b) $g(x)=x^5+x$  \hspace{1cm}   
  c) $h(x)=2x-x^2$ \\

  \textbf{a)} $f(-x)=1-(-x)^4=1+x^4=f(x)$, logo $f$ é par. \\
  \textbf{b)} $g(-x)=(-x)^5-(-x)= -x^5+x=-(x^5-x)=-g(x)$, logo $g$ é ímpar. \\
  \textbf{c)} Nesse caso temos que $h(1)=2.1-1^2=1$ e $h(-1)=2(-1)-(-1)^2==-3$,
  logo $h$ não é nem par, nem ímpar.
\end{exe}

\begin{figure}[h!]
  \centering
  \includegraphics[scale=1.2]{f1}
  \includegraphics[scale=1.2]{f2}
  \includegraphics[scale=1.2]{f3}
  \caption{ Gráfico das funções $f(x)=1-x^4$, $g(x)=x^5+x$ e $h(x)=2x-x^2$,
  respectivamente.}
\end{figure}

%------------------------------------------------------------------------------%
\subsection{Zeros de função}

Os valores de $x$ que satisfazem a equação $f(x)=0$ são chamados \textbf{zeros}
de $f$. Esses valores correspondem aos interceptos-x do gráfico da função.

\begin{exe}
  Encontre os zeros das seguintes funções:
  a) $f(x)=x^3+2x^2-5x$ 
  b) $g(x)= \frac{x^2-4}{x+3}$ \\

  \textbf{a)} Para encontrarmos os zeros de $f$ precisamos resolver a equação
  $x^3+2x^2-5x=0 \Rightarrow x(x^2+2x-5)=0$, portanto $x=0$ ou $x^2+2x-5=0$.
  Para encontrarmos as raízes de $x^2+2x-5=0$ usamos a fórmula de Bhaskara, uma
  vez que se trata de uma equação do segundo grau. Com isso, $x^2+2x-5=0
  \Rightarrow x=-1+\sqrt{6}$ ou $x=-1-\sqrt{6}$. Logo, $f$ possui três zeros, a
  saber $x=0$, $x=-1+\sqrt{6}$ e $-1-\sqrt{6}$. \\
  
  \textbf{b)} Para encontrarmos os zeros de $g$ precisamos resolver a equação
  $\frac{x^2-4}{x+3}=0 \Rightarrow x^2-4=0 \Rightarrow x=2 \hspace{0,3cm} ou
  \hspace{0,3cm} x=-2$. Logo $g$ possui dois zeros, a saber $x=-2$ e $x=2$. 
\end{exe}

%------------------------------------------------------------------------------%
\subsection{Monotonicidade}

\begin{defin} 
  Seja $f$ uma função definida em um intervalo $I$, dizemos que
  \begin{enumerate}
  \item $f$ é \textbf{crescente} em $I$ se, dados quaisquer $x_1$ e $x_2$ em
    $I$, tais que $x_1< x_2$, tivermos $f(x_1)< f(x_2)$;
  
  \item $f$ é \textbf{decrescente} em $I$ se, dados quaisquer $x_1$ e $x_2$
    em $I$, tais que $x_1< x_2$, tivermos $f(x_1)> f(x_2)$;
  
  \item $f$ é \textbf{constante} em $I$ se, dados quaisquer $x_1$ e $x_2$ em
    $I$, tivermos $f(x_1)=f(x_2)$

  \end{enumerate}
\end{defin}

\begin{figure}[h!]
  \centering
  \includegraphics[scale=0.7]{monoti}
  \caption{$f$ é crescente nos intervalos $(a,b)$ e $(d,e)$, decrescente nos
  intervalos $(b,c)$ e $(e,g)$ e constante no intervalo $(c,d)$. Figura
  retirada do livro pré-cálculo do autor Francisco Magalhães Gomes.}
\end{figure}

%------------------------------------------------------------------------------%
\subsection{Função Injetora}

\begin{defin}(\textbf{Função Injetora}): 
  Uma função $f$, definida em seu domínio $D$ é \textbf{injetora} quando dados
  quaisquer valores reais $x_1$, $x_2 \in D$, se $x_1 \neq x_2$, então, $f(x_1)
  \neq f(x_2)$.
\end{defin}

Em outras palavras, se $f$ é uma função injetora então,
\[
  f(x_1)=f(x_2) \Leftrightarrow x_1=x_2.
\]

\textbf{Teste da reta horizontal:} 
Uma função é injetora se e somente se nenhuma reta 
horizontal intercepta seu gráfico mais de uma vez.

\begin{figure}[h!]
  \centering
  \includegraphics[scale=1.4]{injet}
  \caption{Gráfico de uma função que não é injetora, 
  uma vez que $f(x_1)=f(x_2)$ para $x_1 \neq x_2$.}
\end{figure}

\begin{exe} 
  Verifique se as funções abaixo são injetoras \\
  a) $f(x)=x^2$ \hspace{1cm} 
  b) $f(x)=x^3$ \hspace{1cm}
  c) $y= \frac{3}{5x-2}$ \\

  \textbf{a)} $f$ não é injetora pois $f(-1)=f(1)=1$. \\

  \textbf{b)} Suponha que $f(x_1)=f(x_2) \Rightarrow x_1^3=x_2^3 \Rightarrow
  x_1^3= \sqrt[3]{x_2^3} \Rightarrow x_1=x_2$, logo $f$ é injetora. \\
  
  \textbf{c)} Suponha que $f(x_1)=f(x_2) \Rightarrow \frac{3}{5x_1-2}=
  \frac{3}{5x_2-2} \Rightarrow 3(5x_1-2)=3(5x_2-2) \Rightarrow 5x_1-2=5x_2-2
  \Rightarrow x_1=x_2.$ Logo, $y$ é injetora. 
\end{exe}

%------------------------------------------------------------------------------%
\subsection{Algumas funções usuais}

A seguir serão apresentados algumas funções usuais, onde algumas delas serão
estudadas com mais detalhes em aulas posteriores.

\textbf{Exemplo 1 (Função Potência):} A função $f(x)=x^ n$, onde $n$ é um
número natural, é chamada de \textbf{função potência}. A figura  \ref{pot}
mostra o gráfico de algumas funções potências.

\begin{figure}[h!] \label{pot}
  \centering
  \includegraphics[scale=1.6]{pot1}
  \includegraphics[scale=1.6]{pot2}
\end{figure}

\begin{itemize}
  \item O domínio de $f$ é $\mathbb{R}$.

  \item O conjunto imagem é $\mathbb{R}$ se o expoente for ímpar e é
    $[0,\infty]$ se o expoente é par.

  \item Há um único zero em $x=0$.

  \item Quando o expoente é par, $f$ é par, quando o expoente é ímpar, $f$ é
    ímpar.

  \item Quando o expoente é par, f é decrescente para $x<0$, e crescente para
    $x>0$, quando o expoente é ímpar, $f$ é crescente em $(- \infty, \infty)$.

\end{itemize}

\textbf{Exemplo 2 (Função Polinomial):} A função $f(x)=a_nx^ n+
a_{n-1}x^{n-1}+...+ a_2x^2+a_1x+a_0$. $n$ é um número inteiro não- negativo e
$a_1,a_2,...a_n$ são constantes reais.

\begin{itemize}
  \item O domínio de $f$ é $\mathbb{R}$.
  \item Pode ter até $n$ zeros.
\end{itemize}
Os exemplos 3 e 4 ilustram exemplos de funções polinomiais.

\textbf{Exemplo 3: (Função do $1^{\circ}$ grau)}
É uma função do tipo $f(x)=ax+b$, com $a,b \in \mathbb{R}$. O gráfico de uma
função do $1^{\circ}$ grau é uma reta. $a$ é o coeficiente angular da reta e
$b$ é a intersecção com o eixo $y$. Para encontrar a equação de uma reta
precisamos de dois pontos, em particular se encontrarmos os interceptos com os
eixos conseguimos obter a equação da reta.

\begin{figure}[h!]
  \centering
  \includegraphics[scale=0.35]{Linear1}
  \includegraphics[scale=0.35]{linear2}
  \caption{Gráfico de uma função do primeiro grau.}
\end{figure}

\textbf{Exemplo 4: (Função quadrática)} É uma função polinomial de grau 2. ou
seja, tem a forma $f(x)=ax^2+bx+c$, com $a,b,c \in \mathbb{R}$ e $a \neq 0$. O
gráfico desse tipo de função é sempre uma parábola e tem concavidade para cima
se $a>0$ e concavidade para baixo se $a<0$. O gráfico de uma parábola pode ser
construído a partir de translações da parábola $ax^2$. A figura a seguir mostra
o exemplo de duas parábolas.

\begin{figure}[h!]
  \centering
  \includegraphics[scale=0.5]{p1}
  \includegraphics[scale=0.8]{p2}
  \caption{Gráfico das funções $y=x^2+x+1$ e $-2x^2+3x+1$.}
\end{figure}

\textbf{Exemplo 5: (Função raíz)} É uma função do tipo $f(x)=\sqrt[n]{x}$, com
$n$ sendo uma constante natural maior que $1$.

\begin{itemize}
 \item O domínio é $\mathbb{R}$ se o índice $n$ é ímpar e é $[0,\infty)$ se o
   índice é par.

 \item O conjunto imagem é $\mathbb{R}$ se o índice $n$ é ímpar e é
   $[0,\infty)$ se o índice é par.

 \item Há um único zero em $x=0$.

 \item A função é crescente em todo o domínio.

 \item Quando o índice é ímpar, $f$ é ímpar.
\end{itemize}

\begin{figure}[h!]
  \centering
  \includegraphics[scale=0.9]{raiz}
  \includegraphics[scale=0.9]{raiz-cubica}
  \caption{Gráficos das funções $y=\sqrt{x}$ e $\sqrt[3]{x}$, respectivamente}
\end{figure}

\textbf{Exemplo 6: (Função recíproca)}
É uma função do tipo $f(x)=\frac{1}{x^n}$, com $x \neq 0$.
\begin{itemize}
  \item o domínio inclui todos os valores reais, exceto o zero.

  \item o conjunto imagem é $(-\infty,0) \cup (0, \infty)$ se o expoente é
    ímpar e $(0, \infty)$ se o expoente é par.

  \item A função não tem zeros.

  \item Quando o expoente é ímpar, a função é decrescente em todos os pontos do
    domínio. Já quando o expoente é par, $f$ é crescente para $x<0$ e
    decrescente para $x>0$.

  \item Quando o expoente é par, $f$ é par. Quando o expoente é ímpar, $f$ é
    ímpar.

\end{itemize}

\begin{figure}[h!]
  \centering
  \includegraphics[scale=0.7]{hiperbole}
  \caption{Gráfico da função $\frac{1}{x}$.}
\end{figure}

\textbf{Exemplo 7: (Função racional)}  É uma função do tipo
$f(x)=\frac{p(x)}{q(x)}$, onde $p(x),q(x)$ são polinômios e $q(x) \neq 0$. Por
exemplo, a função $\frac{2x^4-x^2+1}{x^2-4}$ é uma função racional e seu
domínio é $\{ x \in \mathbb{R} | x \neq \pm 2 \}$.

\textbf{Exemplo 8: (Função exponencial)} É uma função do tipo $f(x)=a^x$. onde
$a$ é uma constante positiva.

\begin{itemize}
  \item O domínio de $f$ é $\mathbb{R}$ (ou $(- \infty, \infty)$).

  \item A imagem é $(0, \infty)$.

  \item A função não possui zeros;

  \item se $a>1$, a função é crescente em todo o domínio e se $0<a<1$ a função
    é decrescente em todo o domínio.

  \item A função não é par nem ímpar.

\end{itemize}

A figura a seguir mostra o gráfico das funções $y=2^x$ e $y=(0.5)^x$

\begin{figure}[h!]
  \centering
  \includegraphics[scale=0.7]{expone}
  \caption{Gráficos das funções $y=2^x$ e $y=(0.5)^x$. Figura retirada do livro
  Cálculo 1 do autor James Stewart.}
\end{figure}

\textbf{Exemplo 9: (funções trigonométricas)} As funções trigonométricas são
funções angulares obtidas através do auxílio do círculo trigonométrico. As
principais funções trigonométricas são as funções $y= \sin(x)$ e $y=\cos(x)$ e
$y=tg(x)$. Essas funções são períodicas, ou seja,  os valores da função $(f(x)
= y)$ se repetem para determinados valores da variável $x$, ou seja, para cada
período determinado pelos valores de $ x$, iremos obter valores repetidos para
a função.

\begin{itemize}
  \item O domínio de ambas funções $y=\sin(x)$ e $y=\cos(x)$ é $\mathbb{R}$.

  \item A Imagem de ambas funções $y=\sin(x)$ e $y=\cos(x)$ é $[-1,1]$, ou
    seja, $-1 \leq \sin(x) \leq 1$ e $-1 \leq \cos(x) \leq 1$.

  \item A função $y=\sin(x)$ é ímpar e a função $y= \cos(x)$ é par.

\end{itemize}

\begin{figure}[h!]
  \centering
  \includegraphics[scale=0.8]{seno}
  \includegraphics[scale=0.8]{cos}
  \caption{Gráfico da função $y=\sin(x)$ e  $y=\cos(x)$, respectivamente.}
\end{figure}

\textbf{Exemplo 10:(Funções definidas por partes)}: Uma função definida por
partes é uma função que tem mais de uma lei para o seu domínio. 
A função 
\[
  f(x)=\left\{\begin{array}{cc}
    x^2-x \,\, se \,\, x < 1\\
    x-1 \,\, se \,\, x \geq 1
  \end{array}\right.
\]
é um exemplo de função definida por partes.

%------------------------------------------------------------------------------%
\section{Transformação de Funções}

Nessa seção usaremos translações, expansões e reflexões para fazer o esboço de
gráfico de algumas funções a partir de funções já conhecidas (seção anterior).
Abaixo, listamos algumas transformações nas variáveis $x$ e $y$ e como elas
afetam o gráfico de um função $y=f(x)$.

\begin{center}
\begin{tabular}{|l|l|} \hline
$f(x) \mapsto f(-x)$ & Reflexão em torno do eixo $y$ \\ \hline
$f(x) \mapsto -f(x)$ & Reflexão em torno do eixo $x$ \\ \hline
$f(x) \mapsto f(x+a)$ & Translação horizontal  \\
& direita se $a < 0$, esquerda se $a > 0$ \\ \hline
$f(x) \mapsto f(x)+a$ & Translação vertical, \\
& abaixo se $a < 0$, acima se $a > 0$ \\ \hline
$f(x) \mapsto f(ax)$ & Mudança de escala horizontal \\
& contração se $a > 1$, expansão se $0<a < 1$ \\ \hline
$f(x) \mapsto af(x)$ & Mudança de escala vertical, $a>0$ \\
& expansão se $a >1$, contração se $0<a < 1$ \\ \hline
\end{tabular}
\end{center}

\begin{center}
\begin{tabular}{|c|c|c|}
\hline
Original & Reflexão em $y$ & Reflexão em $x$  \\ \hline
\includegraphics[scale=0.25]{geral} &
\includegraphics[scale=0.25]{reflecty} &
\includegraphics[scale=0.25]{reflectx} \\ \hline
$y=f(x)$ & $y=f(-x)$ & $y=-f(x)$ \\ \hline
\end{tabular}
\end{center}

\begin{center}
\begin{tabular}{|c|c|c|c|}
\hline
Translação direita & Translação esquerda & Translação acima & Translação abaixo \\ \hline
\includegraphics[scale=0.25]{transh} &
\includegraphics[scale=0.25]{transh2} &
\includegraphics[scale=0.25]{transv} &
\includegraphics[scale=0.25]{transv2}\\ \hline
$y=f(x+a), a< 0$ & $y=f(x+a), a> 0$ & $y=f(x)+a, a>0$ & $y=f(x)+a, a<0$ \\ \hline
\end{tabular}
\end{center}

\begin{center}
\begin{tabular}{|c|c|c|c|}
\hline
Expansão horizontal & Contração horizontal & Expansão vertical & Contração vertical \\ \hline
\includegraphics[scale=0.25]{scaleh} &
\includegraphics[scale=0.25]{scaleh2} &
\includegraphics[scale=0.25]{scalev} &
\includegraphics[scale=0.25]{scalev2} \\ \hline
$y=f(ax), 0<a< 1$ & $y=f(ax), a> 1$ & $y=af(x), a>1$ & $y=af(x), 0<a<1$  \\ \hline
\end{tabular}
\end{center}

\begin{exe} 
  A partir do gráfico de $f(x)= \sqrt{x}$, 
  use transformações para obter os gráficos de  \\
  a) $y=\sqrt{x}-1$ \hspace{1cm} 
  b) $y=\sqrt{x-2}$ \hspace{1cm} 
  c) $y=- \sqrt{x}$ \hspace{1cm} 
  d) $y=2\sqrt{x}$  \hspace{1cm} 
  e) $y= \sqrt{-x}$
\end{exe}

%------------------------------------------------------------------------------%
\subsection{Função composta}

\begin{defin} \label{com}\textbf{Função composta}] 
  Dadas as funções $f $ e $ g$ cujos domínios são $A$ e $B$, respectivamente,
  definimos a \textbf{função composta} $f \circ g$ por
  \[
    (f \circ g)(x)= f(g(x))
  \]
  O domínio de $f \circ g$ é o conjunto dos valores de $x \in B$ tais que $g(x)
  \in A$
\end{defin}

\begin{exe}
  Se $f(x)= x^2$ e $g(x)=x-3$, encontre as compostas 
  $f \circ g$ e $g \circ f$. \\
  Por definição 
  $f \circ g= f(g(x))=f(x-3)=(x-3)^2$ e 
  $g \circ f= g(f(x))=g(x^2)= x^2-3$.
\end{exe}

\begin{exe} 
  Se $f(x)= \frac{x}{x+1}$, $g(x)=x^{10}$ e $h(x)=x+3$, 
  encontre  $f \circ g \circ h$.
  \[  
    f \circ g \circ h
    = f(g(h(x)))
    = f(g(x+3))
    =f((x+3)^{10})
    =\frac{(x+3)^{10}}{(x+3)^{10}+1}.
  \]
\end{exe}

%------------------------------------------------------------------------------%

\begin{obs} \label{obs} 
  Devemos está atentos ao calcularmos o domínio da função composta. Segundo a
  definição \ref{com} para que sejamos capazes de calcular $f(g(x))$ é preciso
  que:
  \begin{itemize}
    \item[1.] $x$ pertença ao domínio de $g$;
    \item[2.] $g(x)$ (o valor de $g$ em $x$) pertença ao domínio de $f$;
  \end{itemize}
  Para deixar mais clara a definição do domínio da função composta, veremos
  como obtê-lo de forma prática no exemplo a seguir.
\end{obs}

\begin{exe}
  Sejam dadas as funções $f(x)=\frac{1}{x^2-4}$ e $g(x)=\frac{1}{x}$. Obtenha
  $f(g(x))$ e determine o domínio dessa composta. \\

  $f(g(x))= f(\frac{1}{x})=\frac{1}{\frac{1}{x^2}-4}=
  \frac{1}{\frac{1-4x^2}{x^2}}= \frac{x^2}{1-4x^2}$. Pela Observação  \ref{obs}
  para que $f(g(x))$ esteja bem definida temos que ter : \\

  \textbf{1.} $x$ pertença ao domínio de $g$, isto é, $x \neq 0$. \\

  \textbf{2.} $g(x)$ pertença ao domínio de $f$, isto é $1-4x^2 \neq 0
  \Rightarrow x \neq -\frac{1}{2}$ e $x \neq \frac{1}{2}$ . \\

  Logo  $D(f(g(x)))= \{ x \in \mathbb{R} | x \neq 0, x \neq -\frac{1}{2}, x
  \neq \frac{1}{2}\}$
\end{exe}

%------------------------------------------------------------------------------%

%------------------------------------------------------------------------------%
\section{Exercícios Propostos}

\begin{enumerate}
  \item Uma companhia de turismo cobra $R\$ 1200,00$ por um pacote turístico
    individual, mas dá um bom desconto quando os turistas viajam em grupo. Se um
    grupo tem $x$ pessoas, com $x \leq 70$, cada integrante paga apenas
  \[C(x)=1200-12x.\]
  Sabendo que a receita da empresa em cada grupo é dada pelo produto do valor pago por pessoa pelo número de membros do grupo.

\begin{enumerate}
\item escreva a função $R(x)$ que fornece a receita obtida pela empresa de turismo com o grupo de $x$ pessoas
\item trace o gráfico da função para $x \in [10,70]$
\item indique qual o tamanho do grupo fornece a maior receita
\end{enumerate}
\item Calcule
  \[\frac{f(x)-f(a)}{x-a}\]
para as funções e os valores de $a$ fornecidos a seguir. Simplifique os resultados e suponha sempre que os denominadores são diferentes de zero.
\begin{itemize}
\item[a)] $f(x)=3x+4, \hspace{1cm} a=2$
\item[b)] $f(x)=x^2+6, \hspace{1cm} a=4$
\end{itemize}

\item O gráfico da função $f(x)= \sqrt{-x^2+2x+3}$ é mostrado a seguir. Determine:

\begin{center}
  \includegraphics[width=0.5\textwidth]{questao_aula1-3}
\end{center}

\begin{itemize}
\item[a)] Os valores de $f(0),f(0,5)$ e $f(2)$
\item[b)] O domínio de $f$
\item[c)]O conjunto imagem de $f$
\item[d)] Os zeros de $f$
\item[e)] Os intervalos de crescimento e decrescimento
\item[f)] Os pontos de máximo e mínimo absoluto
\end{itemize}

\item Encontre o domínio das seguintes funções: \\
a) $f(x)= \frac{1}{x}+ \frac{1}{x+5}$ \hspace{1cm} b) $f(x)= \sqrt{x+2}+ \sqrt{x-5}$ \hspace{1cm} c) $f(x)=\frac{\sqrt{2x-3}}{x^2-5x+4}$ \hspace{1cm} d) $h(x)= \frac{1}{\sqrt[4]{x^2-5x}}$

\item Verifique algebricamente se as funções a seguir são pares, ímpares ou não possuem simetria. \\
a) $f(x)= x^2-3$ \hspace{1cm} b) $f(x)=-2x$ \hspace{1cm}  c) $f(x)=\frac{1}{x^2+1}$\hspace{1cm} d) $f(x)=2x^5-x^3+x$   \hspace{1cm} e) $y=x \sqrt{x}$ \hspace{1cm} f) $f(x)=2x-1$.

\item  Se $f(x)=x+2$, $g(x)=x^2-4$ e $h(x)=-2x+5$
\begin{enumerate}
\item Verifique se cada uma dessas funções são par, ímpar ou nem par nem ímpar
\item Verifique se essas funções são injetoras
\item Esboce o gráfico dessas funções, encontre o conjunto imagem e os intervalos de crescimento e decrescimento
\item Encontre o valor de  $\frac{g(1)-f(0)-h(3)}{h(1)}$
\item Calcule $g(h(2)))$
\item A partir do gráfico de $y=2^x$, esboce o gráfico de $y= 2^{f(x)}$
\end{enumerate}


\item (\textbf{Cercando um pasto}) Um fazendeiro pretende usar $400 m$ de cerca para delimitar uma área retangular que servirá de pasto. Responda os itens a seguir , lembrando que a área de um terreno retangular com largura $l$ e profundidade $p$ é dada por $l.p$, e que o perímetro desse terreno retangular é igual a $2l+2p$.
\begin{itemize}
\item[a)] Relacione a profundidade com a largura do pasto, considerando o uso dos $400 m$ de cerca.
\item[b)] Escreva uma função que forneça a área cercada em relação a largura do pasto.
\item[c)] Calcule a área do pasto, supondo que sua largura é $75m$. Faça o mesmo para uma largura de $150m$.
\item[d)] Determine o domínio da função área que você obteve.
\item[e)] Esboce o gráfico da função área.
\item[f)] Indique os intervalos nos quais a função é crescente e em quais ela é decrescente.
\item[g)] Com base no gráfico, indique se é possível cercar uma área de $12000 m^2$.
\item[h)] Ainda com base no gráfico, determine a maior área que pode ser cercada e as dimensões do terreno nesse caso.
\end{itemize}

\item A partir do gráfico de $f(x)=x^2$, esboce o gráfico das seguintes funções: \\
a) $g(x)= x^2-3$ \hspace{1cm} b) $h(x)= 3x^2-8$ \hspace{1cm} c)$p(x)=-x^2+6$ \hspace{1cm} d) $q(x)= \frac{1}{2}(x-2)^2$  \\ e) $u(x)=(\frac{1}{4}x)^2$.

\item A partir do gráfico de $f(x)=\sqrt[3]{x}$, esboce o gráfico das seguintes funções: \\
a) $g(x)= \sqrt[3]{x+3}$ \hspace{1cm} b) $h(x)= \sqrt[3]{3x}$ \hspace{1cm} c)$i(x)=-\sqrt[3]{x}$ \hspace{1cm} d) $j(x)= 3\sqrt[3]{x}$  \hspace{1cm} e) $k(x)=1-\sqrt[3]{x-1}$ \hspace{1cm} f)$l(x) =\sqrt[3]{2x}-1$.

\item Dadas as funções $f$ e $g$ a seguir, defina $f(g(x))$, $g(f(x))$, $f(f(x))$ e $g(g(x))$. \\
a) $f(x)=4x$, \hspace{0.1cm} $g(x)= \frac{x^2}{4}$ \hspace{1cm} b )$f(x)=\sqrt{x}$, \hspace{0.1cm} $g(x)= \frac{x}{3}$.

\item Dadas as funções $f(x)=\frac{1}{x-4}$ e $g(x)=x^2$,
\begin{itemize}
\item[a)]Defina $f(g(x))$ e $g(f(x))$ e seus domínios.
\item[b)] Calcule $f(g(-3))$ e $g(f(7))$.

\end{itemize}
\item Na cidade de Salicilina, o número de casos de determinada doença viral varia ao longo do ano, sendo dado aproximadamente por $d(t)=1,9245t^4-53,485t^3+450,18t^2-1081,5t+837,37$, em que $t$ é o mês do ano. Por sua vez, o custo em reais do tratamento dessa doença é dado pela função $c(x)=600x+12000$, em que $x$ é o número de pessoas infectadas.
\begin{itemize}
\item[a)] Determine $c(d(t))$.
\item[b)] Indique o que essa função composta representa.
\item[c)] Calcule o custo aproximado de tratamento dos salicilinenses infectados no mês de abril. (Use uma calculadora).
\end{itemize}
\end{enumerate}

%------------------------------------------------------------------------------%
\subsection{Gabarito:}

\begin{enumerate}

\item
\begin{enumerate}
\item$R(x)=x(1200-12x)$
\item
\begin{center}
  \includegraphics[width=0.5\textwidth]{gabarito_aula1-1b}
\end{center}
\item$x=50$ \\
\end{enumerate}

\item
\begin{enumerate}
\item $3$, se $x \neq 2$
\item $x+4$, se $x \neq 4$ \\
\end{enumerate}

\item
\begin{enumerate}
\item $f(0)=f(2)= \sqrt{3}$ e $f(0,5)=\frac{\sqrt{15}}{2}$
\item $D(f)= \{ x \in \mathbb{R} | -1 \leq x \leq 3\}$
\item Note que o valor mínimo de $f$ é $0$ e o valor máximo é $2$, assim, $Im(f)= [0,2]$
\item $x=-1$ e $x=3$
\item Pelo gráfico podemos notar que $f$ é crescente no intervalo $[-1,1]$ e decrescente no intervalo $[1,3]$
\item $(-1,0)$ e $(3,0)$ são os pontos de mínimo absolutos e $(1,2)$ é o ponto de máximo absoluto
\end{enumerate}



\item
\begin{enumerate}
\item $D(f)=\{ x \in \mathbb{R}| x \neq 0, x \neq -5\}$
\item $D(f)=\{ x \in \mathbb{R}| x \geq 5\}$
\item $D(f)=\{ x \in \mathbb{R}| x \geq \frac{3}{2}, x \neq 4\}$
\item $D(h)=\{ x \in \mathbb{R}| x < 0$ ou $ x > 5\}$   \\
\end{enumerate}

\item
\begin{enumerate}
\item par
\item ímpar
\item par
\item ímpar
\item nem par, nem ímpar
\item nem par, nem ímpar\\
\end{enumerate}

\item
\begin{enumerate}
\item $f$ é nem par, nem ímpar, $g$ é par e $h$ é nem par, nem ímpar
\item $f$ e $h$ são injetoras e $g$ não é injetora
\item
$f$ é sempre crescente, $g$ é decrescente no intevalo $(-\infty,0)$ e crescente no intevalo $(0, \infty)$ e $h$ é sempre decrescente
\begin{figure}[h!]
\centering
\includegraphics[scale=0.55]{gabarito_aula1-6cI}
\caption{Gráfico da função $f(x)=x+2$}
\end{figure}

\begin{figure}[h!]
\centering
\includegraphics[scale=0.55]{gabarito_aula1-6cII}
\caption{Gráfico da função $g(x)=x^2-4$}
\end{figure}

\begin{figure}[h!]
\centering
\includegraphics[scale=0.55]{gabarito_aula1-6cIII}
\caption{Gráfico da função $h(x)=-2x+5$}
\end{figure}
\item $-\frac{4}{3}$
\item $g(h(2))=-3$
\item $ $
  \begin{figure}[h]
  \centering
  \includegraphics[scale=0.55]{gabarito_aula1-6f}
  \caption{Gráfico da função $f(x)=2^{f(x)}$}
  \end{figure}
\end{enumerate}

\item
\begin{enumerate}
\item $D(f)=\{ x \in \mathbb{R}|  x \neq 3\}$
\item $f$ não possui zeros
\item nem par, nem ímpar
\item Intercepta o eixo $y$ em $x=-\frac{1}{3}$ e não intercepta o eixo $x$, visto que não possui zeros
\item $x=-2$ e $x=1$\\
\end{enumerate}

\item
\begin{enumerate}

\item $p=200-l$
\item $A(l)=200l-l^2$
\item $A(75)=9375 m^2$ e $A(150)=7500m^2$
\item $D=(0,200)$, para chegar nesse domínio você deve impor que $A(l)> 0$
\item $ $
        \begin{figure}[h]
        \centering
        \includegraphics[height=6cm , width =10cm]{gabarito_aula1-8e}
        \caption{Gráfico da função $A(l)= 200l-l^2$ }
        \end{figure}
\item É crescente no intervalo $(0,100) $ e decrescente no intervalo $(100,200)$
\item Não, pois a equação $200l-l^2=12000$ não possui solução real
\item $A_{max}=10000m^2$ e as dimensões são $l=p=100$ \\
\end{enumerate}

\item
\begin{enumerate}
    \item $ $
        \begin{figure}[h]
        \centering
        \includegraphics[scale=0.55]{gabarito_aula1-9a}
        \caption{Gráfico da função $g(x)= \sqrt[3]{x+3}$ }
        \end{figure}
    \item $ $
        \begin{figure}[h]
        \centering
        \includegraphics[scale=0.55]{gabarito_aula1-9b}
        \caption{Gráfico da função $h(x)= \sqrt[3]{3x}$ }
        \end{figure}
    \item$ $
        \begin{figure}[h]
        \centering
        \includegraphics[scale=0.55]{gabarito_aula1-9c}
        \caption{Gráfico da função $i(x)= -\sqrt[3]{x}$ }
        \end{figure}
    \item$ $
        \begin{figure}[h]
        \centering
        \includegraphics[scale=0.55]{gabarito_aula1-9d}
        \caption{Gráfico da função $j(x)= 3\sqrt[3]{x}$ }
        \end{figure}
    \item$ $
        \begin{figure}[h]
        \centering
        \includegraphics[scale=0.55]{gabarito_aula1-9e}
        \caption{Gráfico da função $k(x)= 1-\sqrt[3]{x-1}$ }
        \end{figure}
    \item $ $
        \begin{figure}[h]
        \centering
        \includegraphics[scale=0.55]{gabarito_aula1-9f}
        \caption{Gráfico da função $l(x)=\sqrt[3]{2x}-1$ }
        \end{figure}
\end{enumerate}

\item
\begin{enumerate}
\item $f(g(x))=x^2$, $g(f(x))=4x^2$, $f(f(x))=16x$ e $g(g(x))=\frac{x^4}{64}$
\item $f(g(x))=\sqrt{\frac{x}{3}}$, $g(f(x))=\frac{\sqrt{x}}{3}$, $f(f(x))=\sqrt{\sqrt{x}}$ e $g(g(x))=\frac{x}{9}$\\
\end{enumerate}

\item
\begin{enumerate}
\item  $f(g(x))= \frac{1}{x^2-4}$ e $D(f(g))=\{ x \in \mathbb{R}| x \neq -2, x \neq 2 \}$, $g(f(x))=\frac{1}{(x-4)^2}$ e $D(g(f))=\{ x \in \mathbb{R}| x \neq 4\}$ \\
\item  $f(g(-3))= \frac{1}{5}$ e $g(f(7))= \frac{1}{9}$
\end{enumerate}

\item
\begin{enumerate}
    \item $c(d(t)) = 1154,7t^4 - 32091t^3 + 270108t^2 - 648900t + 514422$
    \item O valor interno $d(t)$ calcula o número de casos para um determinado mês ($t$).
    O valor externo $c(d(t))$ calcula o custo total baseado nesse número de casos.
    \item $d(4) \approx783,33$ pessoas. Como o número de pessoas é inteiro devemos calcular o custo para 784 pessoas. Assim $c(784) = \text{R\$ } 482.400,00$
\end{enumerate}
\end{enumerate}



%------------------------------------------------------------------------------%
